\section{Результаты инференса}
\subsection{Результаты инференса случайных примеров из тестового набора данных}

Маски, полученные на выходе модели при обработке случайных примеров из тестового набора данных, полностью перекрывают истинные координаты полюсов/нулей %(рис.~\ref{fig:test_sample1}, \ref{fig:test_sample2}, \ref{fig:test_sample3}, \ref{fig:test_sample4}).

\subsection{Результаты инференса данных эксперимента}

Использовались как характеристики, измеренные внутрисхемно алгоритмом \ref{sec:fr_optim} в приложении \enquote{Digital Points} (приложение~\ref{appendix:application}), так и измеренные сертифицированным измерительным прибором AP300 Frequency Response Analyzer (приложение~\ref{appendix:bench}).

\begin{figure}[H]
	\centering
	\includegraphics [width=1.00\linewidth] {parametersEstimation/G_I_PI.pdf}

	\caption{Результаты эксперимента}
	\label{fig:G_I_PI_detection}
\end{figure}

\begin{figure}[H]
	\centering
	\includegraphics [width=1.00\linewidth] {parametersEstimation/G_U_PI.pdf}

	\caption{Результаты эксперимента}
	\label{fig:G_U_PI_detection}
\end{figure}

\begin{figure}[H]
	\centering
	\includegraphics [width=1.00\linewidth] {parametersEstimation/G_I_PID.pdf}

	\caption{Результаты эксперимента}
	\label{fig:G_I_PID_detection}
\end{figure}

\begin{figure}[H]
	\centering
	\includegraphics [width=1.00\linewidth] {parametersEstimation/Zout.pdf}

	\caption{Результаты эксперимента}
	\label{fig:Zout_detection}
\end{figure}

\begin{figure}[H]
	\centering
	\includegraphics [width=1.00\linewidth] {parametersEstimation/ZoutI.pdf}

	\caption{Результаты эксперимента}
	\label{fig:ZoutI_detection}
\end{figure}

\begin{figure}[H]
	\centering
	\includegraphics [width=1.00\linewidth] {parametersEstimation/Wui.pdf}

	\caption{Результаты эксперимента}
	\label{fig:Wui_detection}
\end{figure}